% ============================================================================
% Fichier  : partie5/P05_C15_secteurs_sensibles.tex
% Titre    : Les Secteurs Sensibles
% Auteur   : MINKA MI NGUIDJOÏ Thierry Emmanuel
% Version  : 1.0 -- novembre 2026
% Corrigés : corriges/Corr_C15.tex
% ============================================================================
% Sources : loi n° 2024/017, art. 4, 5, 9, 18, 19(2), 40(3), 45, 48, 72
% (texte officiel) ; Minka Mi Nguidjoï, Massoussi Songo et Bialo (2026,
% manuscrit) pour le régime spécial de la justice.
% Faits nouveaux (bible) : vidéosurveillance de la salle d'attente de
% pédiatrie de MBOLO ; transmission des résultats aux mutuelles ; audit de
% suivi d'octobre 2026.
% ============================================================================

\chapter{Les Secteurs Sensibles}
\label{chap:secteurs}
\chaptermark{Les secteurs sensibles}

\epigraph{%
    « Quand la loi renvoie à un texte particulier,\\
    l'auditeur cherche le texte avant de conclure. »%
}{Principe du cours}

\niveauChapitre{Expert}{%
    Chapitres~\ref{chap:donnees}, \ref{chap:principes} et~\ref{chap:loi}
}

\begin{boxobjectifs}
\begin{itemize}
    \item Identifier les traitements que la loi renvoie à des textes
          particuliers (articles~4 et~72).
    \item Auditer quand régime général et régime spécial se rencontrent.
    \item Appliquer les règles propres aux données de santé, aux
          transactions bancaires et aux données des mineurs.
    \item Auditer une vidéosurveillance~: fondement, information, durée,
          lieux filmés.
    \item Situer l'expert, l'OPJ et le magistrat au regard de la loi.
\end{itemize}
\end{boxobjectifs}

\bigskip

% ============================================================================
\section*{Avant de commencer : une caméra dans la salle d'attente}
% ============================================================================

Bertoua, mardi 6~octobre~2026. Pour l'audit de suivi du CENTRE MÉDICAL
MBOLO, Mme~BELLO visite la pédiatrie. Une caméra filme la salle d'attente,
où des mères allaitent et où des enfants attendent leurs vaccins. Les images
sont affichées à l'accueil et enregistrées trente jours. Aucun panneau ne
le signale. « C'est pour la sécurité, explique le surveillant, il y a eu des
vols de téléphones. »

Au même moment, le service de facturation transmet chaque mois aux
mutuelles des patients la liste des actes, avec le diagnostic en clair,
« parce que les mutuelles le demandent ». Aucun texte n'a été montré à
l'auditrice qui oblige à transmettre le diagnostic.

\begin{boxalert}
\textbf{Violation typique.} Justifier un traitement de données de santé ou de
mineurs par la seule bonne intention (« la sécurité », « le remboursement »)
sans vérifier le fondement, la proportion et les règles propres au secteur.

\smallskip
\textbf{Règle~:} dans un secteur sensible, chaque traitement cumule les
exigences de la loi générale, les règles propres aux données sensibles et,
quand il existe, le texte particulier du secteur. Ce qui n'est pas prévu
par un texte reste soumis à la loi générale.
\end{boxalert}

% ============================================================================
\section{Les Renvois à des Textes Particuliers}
\label{sec:c15-renvois}
% ============================================================================

Deux articles de la loi renvoient à d'autres textes.

\begin{boxjuris}
\textbf{\loicam{n°~2024/017}, article~4}~\cite{loi2024017}~:
\textit{« Tout traitement des données à caractère personnel effectué par les
autorités compétentes en matière de sécurité et de défense est régi par
des textes particuliers. »}

\medskip
\textbf{\loicam{n°~2024/017}, article~72}~\cite{loi2024017}~:
\textit{« Le traitement des données à caractère personnel relatif notamment
à la sécurité, à la défense, à la santé, à la justice et à l'état civil est
régi par des textes spécifiques, en tant que de besoin. »}

\smallskip
\textbf{Portée.} L'article~4 vise des \textit{autorités}~; l'article~72 des
\textit{domaines}, avec « notamment » et « en tant que de besoin ». Ce
dernier n'exclut pas ces domaines de la loi générale~: il annonce des
textes qui la complèteront ou l'adapteront. Tant qu'un texte spécifique
n'est pas identifié, l'auditeur applique la loi générale et signale la
réserve.
\end{boxjuris}

\begin{boxprocedure}
\textbf{Auditer un traitement d'un secteur visé par l'article~72.}
\begin{enumerate}
    \item Rechercher le texte spécifique au secteur et au traitement~; noter
          la recherche et son résultat.
    \item S'il existe, l'appliquer en même temps que la loi générale, et
          relever les points où il déroge.
    \item S'il n'existe pas, appliquer la loi générale, et coter \NEK{} de
          critère les seules exigences qu'un texte spécifique pourrait
          modifier.
    \item Appliquer dans tous les cas les règles des données sensibles
          (articles~5, 48, 66) et des mineurs (articles~9(3), 18, 19(2),
          40(3)).
    \item Vérifier les secrets professionnels du secteur, qui s'ajoutent à
          la loi.
\end{enumerate}
\end{boxprocedure}

% ============================================================================
\section{La Santé}
\label{sec:c15-sante}
% ============================================================================

Les données de santé sont sensibles (article~5). Le soin trouve un
fondement dans la troisième dérogation de l'article~9~: le traitement
« nécessaire~[\ldots] à la préservation de la santé de la personne
concernée ou d'une autre personne physique ». Ce fondement couvre le soin,
pas tout ce qui l'entoure~: la facturation, la recherche, la formation, la
communication avec les mutuelles ou les employeurs exigent chacun leur
propre fondement.

Les professionnels de santé sont en outre tenus au secret médical, protégé
par le droit pénal et par la déontologie médicale. Le secret ne se partage
qu'entre ceux qui participent à la prise en charge, et dans la mesure
nécessaire~: c'est la règle que l'auditeur applique aux journaux de
consultation (Chapitre~\ref{chap:securite}), et c'est elle qui interdit de
transmettre un diagnostic à une mutuelle sans texte ou sans consentement
du patient.

% ============================================================================
\section{La Banque et la Finance}
\label{sec:c15-banque}
% ============================================================================

Les transactions bancaires sont des données sensibles (article~5), et
l'article~48, alinéa~2, interdit leur traitement « sans autorisation
préalable des administrations et structures compétentes, dans les
conditions et suivant les modalités prévues par les lois en vigueur ». La
loi ne désigne pas ces administrations~; il peut s'agir des autorités de
régulation bancaire ou monétaire. Pour l'auditeur, une microfinance comme
NYANGA ou une fintech comme NOUN PAY doit pouvoir montrer, outre la demande
d'autorisation à l'Autorité de protection, les agréments et autorisations
de son régulateur sectoriel, et le texte qui encadre l'usage des données
dans son activité. Faute de désignation claire, l'exigence « autorisation
des administrations compétentes au titre de l'article~48(2) » est un
\NEK{} de critère.

Le secteur financier connaît aussi des obligations de conservation et de
transmission (lutte contre le blanchiment, déclarations aux autorités) qui
fondent des traitements au titre de l'obligation légale (article~9(2)).
L'auditeur demande le texte précis~; une « obligation réglementaire »
invoquée sans référence reste une déclaration.

% ============================================================================
\section{Les Mineurs}
\label{sec:c15-mineurs}
% ============================================================================

La loi protège les mineurs à quatre endroits~: le consentement d'un mineur
de dix-huit ans n'est valable qu'appuyé par celui de ses parents ou de son
représentant légal (article~9(3))~; les traitements de leurs données en vue
d'une offre de services doivent être « adéquat[s], pertinent[s] et
limité[s] » (article~18)~; l'interconnexion de fichiers de données
sensibles relatives aux mineurs exige une autorisation propre
(article~19(2))~; les parents peuvent s'opposer au traitement des données
de leur enfant collectées sans leur consentement (article~40(3)).

En pédiatrie, l'enfant est soigné sur le fondement de la préservation de sa
santé, et les parents consentent aux actes~; le consentement aux autres
traitements (photographies pour un dossier de formation, participation à
une étude) doit être recueilli auprès des parents, et l'enfant informé à sa
mesure.

% ============================================================================
\section{La Vidéosurveillance}
\label{sec:c15-video}
% ============================================================================

La vidéosurveillance d'un lieu ouvert au public filme des personnes
identifiables~: c'est un traitement de données personnelles. La loi n'en
fait pas un régime particulier~; elle s'applique donc entièrement, et la
difficulté vient du fondement (Chapitre~\ref{chap:principes})~: ni
l'intérêt légitime ni le contrat n'existent dans la loi, et le
consentement de chaque personne filmée est irréaliste. L'auditeur cote
\NEK{} le fondement tant qu'aucune autorisation ne l'a précisé, et vérifie
tout ce qui ne dépend pas du fondement~: la finalité, les lieux filmés,
l'information par panneau, la durée de conservation, les personnes qui
voient les images.

\begin{tableaucours}{L{5.2cm}L{6.8cm}}{%
    \tableauHeader{Question} &
    \tableauHeader{Ce que l'auditeur vérifie}
}
Finalité & Une seule, écrite~: sécurité des biens et des personnes~; pas de
contrôle du travail \\
Lieux filmés & Pas de lieu de soin, de toilettes, de vestiaire, d'espace
d'allaitement \\
Information & Panneau visible à l'entrée de la zone filmée \\
Durée & Courte, fixée, appliquée (quelques semaines au plus, sauf incident) \\
Accès aux images & Personnes désignées, consultation journalisée \\
Affichage & Écrans hors de la vue du public \\
\end{tableaucours}
\vspace{-0.8em}
\begin{center}{\small\itshape Tableau~15.1. Les questions d'audit d'une vidéosurveillance}\end{center}

% ============================================================================
\section{La Justice et la Chaîne Pénale}
\label{sec:c15-justice}
% ============================================================================

La justice figure parmi les domaines de l'article~72. L'expert judiciaire,
l'officier de police judiciaire et le magistrat traitent tous des données
personnelles~: saisies, analyses, rapports, transmissions. Leur qualité au
regard de la loi (responsable, sous-traitant, autre acteur habilité) n'est
pas fixée par un texte.

Dans leur travail en cours, l'auteur de ce cours et ses
coauteurs~\cite{minka2026chainepenale} proposent de déterminer cette qualité
« opération par opération », et non selon la qualité institutionnelle~:
un expert qui collecte des données sur instruction d'une autorité
judiciaire peut être sous-traitant pour cette opération, et devenir
responsable pour la copie qu'il conserve pour ses propres besoins
techniques. Le Chapitre~\ref{chap:recherche} présente cette grille comme
une piste de recherche. Pour l'audit d'un cabinet d'expertise, elle
suggère une règle pratique~: tenir un registre de ses traitements, et
pour chaque mission, documenter qui a déterminé la finalité.

% ============================================================================
\section{L'Audit de Suivi de MBOLO}
\label{sec:c15-mbolo}
% ============================================================================

\begin{boxscenario}{Mission MBOLO. Les traitements sensibles, octobre 2026}
\textbf{Constatations.} La caméra de pédiatrie filme un espace où des mères
allaitent~; aucun panneau~; écran visible depuis le couloir. Le diagnostic
est transmis en clair aux mutuelles chaque mois, pour 1\,640 patients en
septembre~; aucun texte n'a été présenté, et les formulaires d'adhésion des
mutuelles ne mentionnent pas cette transmission. Les consentements
parentaux sont recueillis pour les actes, pas pour les photographies prises
pour le dossier de formation des internes. Le texte spécifique au secteur
de la santé annoncé par l'article~72 n'a pas été trouvé.
\end{boxscenario}

\begin{boxlecture}{Traitements sensibles de MBOLO, octobre 2026}
\begin{tableaucours}{L{7.8cm}cL{3.2cm}}{%
    \tableauHeader{Exigence} &
    \tableauHeader{Cote} &
    \tableauHeader{Preuve}
}
Le soin repose sur la dérogation de l'article~9(2) & \OUI{} & Finalité de
soin \\
La transmission du diagnostic aux mutuelles a un fondement & \NON{} &
Aucun texte, aucun consentement \\
Les lieux filmés excluent les espaces intimes & \NON{} & Salle
d'allaitement filmée \\
La vidéosurveillance est signalée & \NON{} & Aucun panneau \\
La vidéosurveillance a un fondement valable & \NEK{} & Fondement incertain \\
Les photographies d'enfants ont le consentement des parents & \NON{} &
Dossier de formation \\
Un texte spécifique à la santé modifie ces exigences & \NEK{} & Texte non
trouvé \\
\end{tableaucours}

\medskip
\textbf{Décompte.} Un \OUI{}, quatre \NON{}, deux \NEK{} sur sept.
\hfill\lecture{0{,}14}{0{,}29}{1{,}38}

\medskip
\textbf{Lecture.} Domaine opaque, mais les quatre \NON{} ne dépendent
d'aucun texte spécifique~: un texte de santé pourrait préciser la
transmission aux mutuelles, il ne rendrait pas licite une caméra dans un
espace d'allaitement.
\end{boxlecture}

% ============================================================================
\section{Ce que les Secteurs demandent à chaque Métier}
\label{sec:c15-metiers}
% ============================================================================

\begin{boxmetier}{Juriste conformité}
Tenir une veille des textes spécifiques annoncés par l'article~72, et une
liste, par traitement, des textes sectoriels qui l'encadrent.
\end{boxmetier}

\begin{boxmetier}{Auditeur de la protection des données}
Documenter la recherche du texte spécifique, appliquer la loi générale en
son absence, et ne coter \NEK{} de critère que les exigences qu'un tel
texte pourrait réellement modifier.
\end{boxmetier}

\begin{boxmetier}{Responsable de traitement}
Faire déplacer la caméra et retirer le diagnostic des listes transmises
aux mutuelles dès le rapport~: ces écarts ne dépendent d'aucun texte à
venir.
\end{boxmetier}

\begin{boxresume}
\begin{itemize}
    \item L'article~4 renvoie la sécurité et la défense à des textes
          particuliers~; l'article~72 annonce des textes spécifiques pour la
          santé, la justice, l'état civil et d'autres domaines.
    \item En l'absence de texte spécifique, la loi générale s'applique, et
          l'auditeur signale la réserve.
    \item La préservation de la santé fonde le soin, pas la facturation, la
          formation ou la transmission aux tiers.
    \item Les mineurs bénéficient de quatre protections propres~; la
          vidéosurveillance se vérifie lieu par lieu.
    \item À MBOLO, quatre écarts sont établis indépendamment de tout texte
          sectoriel.
\end{itemize}
\end{boxresume}

\begin{boxinfo}
\textbf{Points de contrôle.} Vous avez compris ce chapitre si vous savez~:
\begin{itemize}
    \item lire l'article~72 sans en faire une exclusion générale~;
    \item auditer une vidéosurveillance~;
    \item dire ce que couvre la dérogation de santé de l'article~9.
\end{itemize}
\end{boxinfo}

\begin{boxexo}[title={\ding{45}\quad Exercice~15.1. Quel fondement \niveaubase}]
Quel fondement pour~: (a) le dossier médical d'un patient~; (b) la
facturation d'un acte~; (c) une étude sur le paludisme à partir des
dossiers~; (d) la photographie d'une plaie pour la formation des
internes~?
\end{boxexo}

\begin{boxexo}[title={\ding{45}\quad Exercice~15.2. Les mutuelles \niveaumoyen}]
Proposez un circuit de transmission aux mutuelles conforme à la loi et au
secret médical~: données transmises, fondement, information des patients.
\end{boxexo}

\begin{boxexo}[title={\ding{45}\quad Exercice~15.3. La caméra \niveaumoyen}]
Rédigez la recommandation sur la vidéosurveillance de la pédiatrie, en
reprenant les questions du Tableau~15.1. Que reste-t-il \NEK{} après sa mise
en œuvre, et pourquoi~?
\end{boxexo}

\begin{boxexo}[title={\ding{45}\quad Exercice~15.4. Un cabinet d'expertise \niveauexpert}]
Un cabinet d'expertise judiciaire vous demande un audit. Pour trois de ses
opérations (recevoir un scellé, analyser une copie, conserver une copie
après la mission), proposez la qualification de l'expert au regard de la
loi, en suivant la grille opération par opération, et dites quelles
exigences vous audit\-erez pour chacune.
\end{boxexo}

\bigskip

\begin{boxinfo}
\textbf{Pour aller plus loin.}
\begin{itemize}
    \item La loi n°~2024/017~\cite{loi2024017}, articles~4, 9, 18, 48 et~72.
    \item T.~E.~Minka Mi Nguidjoï, R.~T.~Massoussi Songo et
          D.~Bialo~\cite{minka2026chainepenale}, pour la chaîne pénale.
\end{itemize}
\end{boxinfo}

\begin{boxinfo}
\textbf{Transition.} Tout au long du cours, des \NEK{} se sont accumulés~: de
preuve et de critère, de faits et de textes. Peut-on mesurer cette
incertitude, et que dit-elle de l'état du droit lui-même~? C'est la
question de recherche du Chapitre~\ref{chap:recherche}.
\end{boxinfo}
